\documentclass[11pt]{article}
\usepackage[]{authblk}
\usepackage{graphicx}
\usepackage{color}
\usepackage{longtable}
\usepackage{hanging}
\usepackage{indentfirst}
\usepackage{setspace}
\usepackage{enumitem}
\usepackage{verbatim}
\usepackage{upgreek}
\usepackage{framed}
\usepackage{ textcomp }
\usepackage{url}
\usepackage{soul}
\usepackage{amsmath, amsfonts,amssymb,mathrsfs}
\usepackage{fancyhdr}
\usepackage[compact]{titlesec}
\usepackage[T1]{fontenc}
\usepackage{lmodern}

\usepackage[backend=bibtex,hyperref=true,citestyle=authoryear,bibstyle=authortitle,firstinits=true,terseinits=true,doi=false,url=false,maxbibnames=10,maxcitenames=2]{biblatex}
\addbibresource{../bib_tex/master_refs.bib}
\DeclareCiteCommand{\cite}
  {\usebibmacro{prenote}}
  {\usebibmacro{citeindex}%
   \printtext[bibhyperref]{\usebibmacro{cite}}}
  {\multicitedelim}
  {\usebibmacro{postnote}}

\DeclareCiteCommand*{\cite}
  {\usebibmacro{prenote}}
  {\usebibmacro{citeindex}%
   \printtext[bibhyperref]{\usebibmacro{citeyear}}}
  {\multicitedelim}
  {\usebibmacro{postnote}}

\DeclareCiteCommand{\parencite}[\mkbibparens]
  {\usebibmacro{prenote}}
  {\usebibmacro{citeindex}%
    \printtext[bibhyperref]{\usebibmacro{cite}}}
  {\multicitedelim}
  {\usebibmacro{postnote}}

\DeclareCiteCommand*{\parencite}[\mkbibparens]
  {\usebibmacro{prenote}}
  {\usebibmacro{citeindex}%
    \printtext[bibhyperref]{\usebibmacro{citeyear}}}
  {\multicitedelim}
  {\usebibmacro{postnote}}

\DeclareCiteCommand{\footcite}[\mkbibfootnote]
  {\usebibmacro{prenote}}
  {\usebibmacro{citeindex}%
  \printtext[bibhyperref]{ \usebibmacro{cite}}}
  {\multicitedelim}
  {\usebibmacro{postnote}}

\DeclareCiteCommand{\footcitetext}[\mkbibfootnotetext]
  {\usebibmacro{prenote}}
  {\usebibmacro{citeindex}%
   \printtext[bibhyperref]{\usebibmacro{cite}}}
  {\multicitedelim}
  {\usebibmacro{postnote}}

\DeclareCiteCommand{\textcite}
  {\boolfalse{cbx:parens}}
  {\usebibmacro{citeindex}%
   \printtext[bibhyperref]{\usebibmacro{textcite}}}
  {\ifbool{cbx:parens}
     {\bibcloseparen\global\boolfalse{cbx:parens}}
     {}%
   \multicitedelim}
  {\usebibmacro{textcite:postnote}}

\newcommand{\citep}{\parencite}
\newcommand{\citet}{\textcite}
\defbibheading{relevref}[\refname]{\section*{Relevant References}}

\renewcommand{\postnotedelim}{\iffieldpages{postnote}{\addcolon}{\addcomma\space}} 
\DeclareFieldFormat{postnote}{#1} 

\DeclareFieldFormat[article, inbook, incollection, inproceedings, patent, thesis, unpublished]{title}{#1}
\DeclareFieldFormat[article, inbook, incollection, inproceedings, patent, thesis, unpublished]{journaltitle}{\mkbibemph{#1}\nopunct}
\DeclareFieldFormat[article, inbook, incollection, inproceedings, patent, thesis, unpublished]{volume}{{#1}\addcolon} %puts volume number in parens
%\DeclareFieldFormat[article, inbook, incollection, inproceedings, patent, thesis, unpublished]{year}{\mkbibparens{#1}\nopunct} %puts year in parens

\DeclareFieldFormat[article, incollection, patent, thesis, unpublished]{pages}{{\nopp#1}}

\DeclareFieldFormat{sentencecase}{\MakeSentenceCase{#1}}

\renewbibmacro*{title}{%
  \ifthenelse{\iffieldundef{title}\AND\iffieldundef{subtitle}}
    {}
    {\ifthenelse{\ifentrytype{article}\OR\ifentrytype{inbook}%
      \OR\ifentrytype{incollection}\OR\ifentrytype{inproceedings}%
      \OR\ifentrytype{inreference}}
      {\printtext[title]{%
        \printfield[sentencecase]{title}%
        \setunit{\subtitlepunct}%
        \printfield[sentencecase]{subtitle}}}%
      {\printtext[title]{%
        \printfield[titlecase]{title}%
        \setunit{\subtitlepunct}%
        \printfield[titlecase]{subtitle}}}%
     \newunit}%
  \printfield{titleaddon}}

\DefineBibliographyStrings{english}{% various adjustments to common bib entry strings
urlseen = {Accessed:},% What goes in front of the date a URL was accessed/retrieved etc.
editor = {(Ed)},%Ed – no dot, in brackets
editors = {(Eds)},% Eds – no dot, in brackets
byeditor = {(Ed.)}}% ‘Edited by’ for edited works

\DeclareNameAlias{default}{last-first}

\renewbibmacro{in:}{}

\renewbibmacro{publisher+location+date}{
  \iflistundef{publisher}
    {}
    {\printlist{publisher}%
       {\addcomma\space}%
      \iflistundef{location}
        {}
        {\printlist{location}}%
    }
}

\DeclareBibliographyDriver{article}{%
\usebibmacro{bibindex}%
\usebibmacro{begentry}%
\usebibmacro{author/translator+others}%
\newunit\newblock
\printfield{year}%
\setunit{\labelnamepunct}\newblock
\usebibmacro{title}%
\newunit
\printlist{language}%
\newunit\newblock
\usebibmacro{byauthor}%
\newunit\newblock
\usebibmacro{bytranslator+others}%
\newunit\newblock
\printfield{version}%
\newunit\newblock
%\usebibmacro{in:}% %mit in:
\usebibmacro{journal}%
\newunit\newblock
\printfield{volume}%
\newunit\newblock
\usebibmacro{byeditor+others}%
\newunit\newblock
\usebibmacro{note+pages}%
\newunit\newblock
\iftoggle{bbx:isbn}
{}%
\newunit\newblock
\usebibmacro{doi+eprint+url}%
\newunit\newblock
\usebibmacro{addendum+pubstate}%
\newunit\newblock
\usebibmacro{pageref}%
\usebibmacro{finentry}}

\DeclareBibliographyDriver{inproceedings}{%
\usebibmacro{bibindex}%
\usebibmacro{begentry}%
\usebibmacro{author/translator+others}%
\newunit\newblock
\printfield{year}%
\setunit{\labelnamepunct}\newblock
\usebibmacro{title}%
\newunit
\printlist{language}%
\newunit\newblock
\usebibmacro{byauthor}%
\newunit\newblock
\usebibmacro{bytranslator+others}%
\newunit\newblock
\printfield{version}%
\newunit\newblock
%\usebibmacro{in:}% %mit in:
\usebibmacro{booktitle}%
\newunit\newblock
\printfield{volume}%
\newunit\newblock
\usebibmacro{byeditor+others}%
\newunit\newblock
\usebibmacro{publisher+location+date}%
\newunit\newblock
\usebibmacro{note+pages}%
\newunit\newblock
\usebibmacro{pageref}%
\usebibmacro{finentry}}

\DeclareBibliographyDriver{book}{%
\usebibmacro{bibindex}%
\usebibmacro{begentry}%
\usebibmacro{author/translator+others}%
\newunit\newblock
\printfield{year}%
\setunit{\labelnamepunct}\newblock
\usebibmacro{title}%
\newunit
\printlist{language}%
\newunit\newblock
\usebibmacro{byauthor}%
\newunit\newblock
\usebibmacro{bytranslator+others}%
\newunit\newblock
%\usebibmacro{in:}% %mit in:
\usebibmacro{booktitle}%
\newunit\newblock
\printfield{volume}%
\newunit\newblock
\usebibmacro{publisher+location+date}%
\newunit\newblock
\usebibmacro{note+pages}%
\newunit\newblock
\usebibmacro{pageref}%
\usebibmacro{finentry}}





\setlength{\evensidemargin}{0in}
\setlength{\headheight}{0in}
\setlength{\headsep}{0in}
\setlength{\oddsidemargin}{-0.25in}
\setlength{\paperheight}{11in}
\setlength{\paperwidth}{8.5in}
\setlength{\tabcolsep}{0in}
\setlength{\textheight}{9in}
\setlength{\textwidth}{7in}
\setlength{\topmargin}{0in}
\setlength{\topskip}{0in}
\setlength{\voffset}{0in}
\parskip = 0.15in
\pagestyle{plain}
\setlength{\parindent}{0cm}

\definecolor{citescol}{RGB}{194,101,1}
\definecolor{urlscol}{RGB}{0,150,206}
\definecolor{linkscol}{RGB}{149,0,207}
\definecolor{mycol}{RGB}{25,23,191}
\definecolor{outputcol}{RGB}{34,139,34}
\definecolor{tcol}{RGB}{165,0,14}


\DeclareMathAlphabet{\msfsl}{T1}{cmr}{m}{it}
\DeclareMathAlphabet{\msyf}{OMX}{pcr}{m}{it}
\newcommand{\alf}{\upalpha}
\newcommand{\hilight}[1]{\colorbox{yellow}{#1}}

\newcommand{\levelone}[1]{
\bigskip
\noindent{\LARGE{\textsc{#1}}}
\vspace {0.05in}
}

\newcommand{\leveltwo}[1]{
\bigskip
\noindent{\Large{\textit{#1}}}
\vspace {-1mm}
}

\newcommand{\descriptionhead}[1]{
\noindent{\textbf{\textit{#1}}}\\ \vspace{-7mm}
}

\newcommand{\dhead}[1]{
\noindent{\textbf{\textit{#1 --}}}
}


\newcommand{\exs}[1]{
\vspace{-4mm}
\begin{itemize}
\item #1 \\ \vspace{-8mm}
\end{itemize}
}

\newcommand{\nbo}[1]{{\color{red}{#1}}}


\newcommand{\stepbullet}{\noindent \textbullet \ }
\newcommand{\mi}[1]{\textbf{\textit{#1}}}


\newcommand{\levelthree}[1]{\textit{#1 --}}


%\bibliographystyle{apalike}
%\bibpunct[; ]{(}{)}{;}{a}{,}{;}


\usepackage[breaklinks]{hyperref}
\usepackage[all]{hypcap}
\hypersetup{colorlinks=true,linkcolor=linkscol,citecolor=citescol,urlcolor=urlscol}


\newcommand{\R}{\texttt{R} }
\newcommand{\TESS}{\texttt{TESS}}
\newcommand{\PBD}{\texttt{PBD}}
\newcommand{\DDD}{\texttt{DDD}}
\newcommand{\Laser}{\texttt{laser}}
\newcommand{\TreePar}{\texttt{TreePar}}
\newcommand{\diversitree}{\texttt{diversitree}}
\newcommand{\RevBayes}{\texttt{RevBayes}}
\newcommand{\Rev}{\texttt{Rev}}
\newcommand{\MrBayes}{\texttt{MrBayes}}
\newcommand{\BEAST}{\texttt{BEAST}}
\newcommand{\PhyloBayes}{\texttt{PhyloBayes}}
\newcommand{\PAML}{\texttt{PAML}}

\let\otheriint\iint
\let\iint\relax
\usepackage{ wasysym }

\usepackage{framed}
\usepackage[]{listings}
%\usepackage{fontspec}
\usepackage{placeins}
\usepackage{epstopdf}

\lstset{breaklines=true}

\definecolor{shadecolor}{RGB}{238,224,229}

\setlength{\tabcolsep}{5pt}
\setlength{\topmargin}{-0.4in}
\setlength{\headheight}{14.5pt}
\pagestyle{fancy}

\newcommand{\taha}[1]{{\textcolor{red}{[TAH comment: #1]}}} % TAH comment

\titlespacing{\section}{0pt}{*0}{*0}
\titlespacing{\subsection}{0pt}{*0}{*0}
\titlespacing{\subsubsection}{0pt}{*0}{*0}

\titleformat{\section}
  {\normalfont\Large\bfseries\color{mycol}}
  {\thesection}{1em}{}

\titleformat{\subsection}
  {\normalfont\large\bfseries\color{mycol}}
  {\thesubsection}{1em}{}

\titleformat{\subsubsection}
  {\normalfont\bfseries\color{mycol}}
  {\thesubsubsection}{1em}{}

% command for MrBayes command-line step
\newcommand{\cl}[1]{{\texttt{\textbf{#1}}}}

\newcommand{\colx}[1]{{\textcolor{tcol}{#1}}}

\newcommand{\mbcl}[1]{\exs{\cl{MrBayes > {#1}}}}

\newcommand{\rbprmt}{RevBayes > } 
\newcommand{\rbcl}[1]{\exs{\cl{\rbprmt{#1}}}}
\newcommand{\rbout}[1]{\exs{\cl{\textcolor{outputcol}{#1}}}}
\newcommand{\rbdn}{{\Large \symbol{126}}} % This makes a copy/pasteable tilde
\newcommand{\rbclml}[1]{\exs{\cl{\ \ \ \ \ \ \ \ \ \ \ {#1}}}}

% text box settings
% requires compiling w/ XeLaTeX
%\newfontfamily\listingsfont[Scale=1.0]{Courier New}
%\lstset{basicstyle=\listingsfont, columns=texcl}
%\defaultfontfeatures{Mapping=tex-text}


\makeatletter
\lst@CCPutMacro\lst@ProcessOther {"2D}{\lst@ttfamily{-{}}{-{}}}
\@empty\z@\@empty
\makeatother


\usepackage{framed}
\usepackage[]{listings}
\usepackage{fontspec}
\usepackage{placeins}
\usepackage{epstopdf}

\lstset{breaklines=true}

\definecolor{shadecolor}{RGB}{238,224,229}


\setlength{\topmargin}{-0.4in}
\setlength{\headheight}{14.5pt}
\pagestyle{fancy}

\newcommand{\taha}[1]{{\textcolor{red}{[TAH comment: #1]}}} % TAH comment

\titlespacing{\section}{0pt}{*0}{*0}
\titlespacing{\subsection}{0pt}{*0}{*0}
\titlespacing{\subsubsection}{0pt}{*0}{*0}

\titleformat{\section}
  {\normalfont\Large\bfseries\color{mycol}}
  {\thesection}{1em}{}

\titleformat{\subsection}
  {\normalfont\large\bfseries\color{mycol}}
  {\thesubsection}{1em}{}

\titleformat{\subsubsection}
  {\normalfont\bfseries\color{mycol}}
  {\thesubsection}{1em}{}

% command for MrBayes command-line step
\newcommand{\cl}[1]{{\texttt{\textbf{#1}}}}

\newcommand{\colx}[1]{{\textcolor{tcol}{#1}}}

\newcommand{\mbcl}[1]{\exs{\cl{MrBayes > {#1}}}}

\newcommand{\rbprmt}{RevBayes > } 
\newcommand{\rbcl}[1]{\exs{\cl{\rbprmt{#1}}}}
\newcommand{\rbout}[1]{\exs{\cl{\textcolor{outputcol}{#1}}}}
\newcommand{\rbdn}{{\Large \symbol{126}}} % This makes a copy/pasteable tilde
\newcommand{\rbclml}[1]{\exs{\cl{\ \ \ \ \ \ \ \ \ \ \ {#1}}}}

% text box settings
% requires compiling w/ XeLaTeX
\newfontfamily\listingsfont[Scale=1.0]{Courier New}
\lstset{basicstyle=\listingsfont, columns=texcl}
\defaultfontfeatures{Mapping=tex-text}


\makeatletter
\lst@CCPutMacro\lst@ProcessOther {"2D}{\lst@ttfamily{-{}}{-{}}}
\@empty\z@\@empty
\makeatother

\begin{document}
\renewcommand{\headrulewidth}{0.5pt}
\headsep = 20pt
\lhead{ }
\rhead{\textsc {RevBayes Tutorial --- Model Selection \& Partitioning}}

\thispagestyle{plain}
\begin{center}

\textbf{\LARGE Phylogenetic Inference using RevBayes}\\\vspace{2mm}
\textbf{\it{\Large Model Selection \& Data Partitioning}}\\\vspace{2mm}
\end{center}

\section*{Overview}
 

This tutorial demonstrates how to set up and perform an analysis that calculates Bayes factors to select among different partition configurations of aligned DNA sequences. 
After selecting the model that is best supported by the data, the exercise continues with basic inference of an unrooted tree topology and branch lengths using Markov chain Monte Carlo (MCMC). 





%
%\subsection*{Analysis Functions}
%
\newpage
\FloatBarrier
\section{Exercise: Model Selection \& Partitioning using Bayes Factors}

\subsection{Introduction}
%
%
%\bigskip
%\section{Model Selection \& Partitioning using Bayes Factors}

Variation in the evolutionary process across the sites of nucleotide 
sequence alignments is well established, and is an increasingly pervasive feature of datasets 
composed of gene regions sampled from multiple loci and/or different genomes.
Inference of phylogeny from these data demands that we adequately model the underlying process heterogeneity; 
failure to do so can lead to biased estimates of phylogeny and other parameters \citep{brown07}.
To accommodate process heterogeneity within and/or between various gene(omic)
regions, we will evaluate the support for various partition schemes 
using Bayes factors to compare the marginal likelihoods of the candidate partition schemes.

Accounting for process heterogeneity involves adopting a `mixed-model' approach, \citep{ronquist03} in which the sequence alignment is first parsed into a number of partitions that are intended to capture plausible process heterogeneity within the data.
The determination of the partitioning scheme is guided by biological considerations regarding the dataset at hand.
For example, we might wish to evaluate possible variation in the evolutionary process within a single gene region ({\it e.g.}, between stem and loop regions of ribosomal sequences), or among gene regions in a concatenated alignment ({\it e.g.}, comprising multiple nuclear loci and/or gene regions sampled from different genomes).
The choice of partitioning scheme is up to the investigator and many possible partitions might be considered for a typical dataset.

Next, a substitution model is specified for each predefined process partition (using a given model-selection criterion, such as Bayes factors).
In this exercise, we assume that each partition evolved under an independent general-time reversible model with gamma-distributed rates across sites (GTR+$\Gamma$). 
Under this model the observed data are conditionally dependent on the exchangeability rates ($\theta$), stationary base frequencies ($\pi$), and the degree of gamma-distributed among-site rate variation ($\alpha$), as well as the unrooted tree topology ($\Psi$) and branch lengths ($\nu$).
We show the graphical model representation of the GTR+$\Gamma$ mode in Figure \ref{gtrgmfig}. 
When we assume different GTR+$\Gamma$ models for each partitions, this results in a composite model, in which all sites are assumed to share a common, unrooted tree topology and proportional branch lengths, but subsets of sites (`data partitions') are assumed to have independent substitution model parameters.
This composite model is referred to as a \emph{mixed model}. %\taha{need to change the graphical model above to have the same symbols as in this paragraph since this refers to the figure \ref{pipeline} and it may be harder to change that figure.}



Finally, we perform a separate MCMC simulation to approximate the joint posterior probability density of the phylogeny and other parameters.  Note that, in this approach, the mixed model is a fixed assumption of the inference ({\it i.e.}, the parameter estimates are conditioned on the specified mixed model), and the parameters for each process partition are independently estimated.

For most sequence alignments, several (possibly many) partition schemes of varying complexity are plausible {\it a priori}, which therefore requires a way to objectively identify the partition scheme that balances estimation bias and error variance associated with under- and over-parameterized mixed models, respectively.
Increasingly, mixed-model selection is based on \textit{Bayes factors} \citep[{\it e.g.},][]{suchard01}, which involves first calculating the marginal likelihood under each candidate partition scheme and then comparing the ratio of the marginal likelihoods for the set of candidate partition schemes \citep{brandley05,nylander04,mcguire07}.
The analysis pipeline that we will use in this tutorial is depicted in Figure \ref{pipeline}.
\begin{figure}[h!]
\centering
\fbox{\includegraphics[width=6.8in,angle=0]{figures/pipeline.eps}}
\caption{\small The analysis pipeline for Exercise 1. We will explore three partition schemes for the conifer dataset.
The first model (the `uniform model', $M_0$) assumes that all sites evolved under a common GTR+$\Gamma$ substitution model.
The second model (the `moderately partitioned' model, $M_1$) invokes two data partitions corresponding to the two gene regions (atpB and rbcL), and assumes each subset of sites evolved under an independent GTR+$\Gamma$ model.
The final mixed model (the `highly partitioned' model, $M_2$) invokes four data partitions---the first partition corresponds to the atpB gene region, and the remaining partitions correspond to the three codon positons of the rbcL gene region---and each data partition is assumed evolved under an independent GTR+$\Gamma$ substitution model.
Note that we assume that all sites share a common tree topology, $\Psi$, and branch-length proportions, $\nu$, for each of the candidate partition schemes.
We perform two separate sets of analyses for each mixed model---a Metropolis-coupled MCMC simulation to approximate the joint posterior probability density of the mixed-model parameters, and a `stepping-stone' MCMC simulation to approximate the marginal likelihood for each mixed model.
The resulting marginal-likelihood estimates are then evaluated using Bayes factors to assess the fit of the data to the three candidate mixed models.  
}
\label{pipeline}
\end{figure}
 
Given two models, $M_0$ and $M_1$, the Bayes factor comparison assessing the relative plausibility of each model as an explanation of the data, $BF(M_0,M_1)$, is:
$$BF(M_0,M_1) = \frac{\mbox{posterior odds}}{\mbox{prior odds}}.$$
The posterior odds is the posterior probability of $M_0$ given the data, $\mathbf X$, divided by the posterior odds of $M_1$ given the data:
$$\mbox{posterior odds} = \frac{\mathbb{P}(M_0 \mid \mathbf X)}{\mathbb{P}(M_1 \mid \mathbf X)},$$
and the prior odds is the prior probability of $M_0$ divided by the prior probability of $M_1$:
$$\mbox{prior odds} = \frac{\mathbb{P}(M_0)}{\mathbb{P}(M_1)}.$$
Thus, the Bayes factor measures the degree to which the data alter our belief regarding the support for $M_0$ relative to $M_1$ \citep{lavine99}:
\begin{align}\label{BFeq1}
BF(M_0,M_1) = \frac{\mathbb{P}(M_0 \mid \mathbf X, \theta_0)}{\mathbb{P}(M_1 \mid \mathbf X, \theta_1)} \div \frac{\mathbb{P}(M_0)}{\mathbb{P}(M_1)}. 
\end{align}
This, somewhat vague, definition does not lead to clear-cut identification of the ``best'' model. Instead, \textsl{you} must decide the degree of your belief in $M_0$ relative to $M_1$. 
Despite the absence of any strict ``rule-of-thumb'', you can refer to the scale \citep[outlined by][]{jeffreys61} for interpreting these measures (Table \ref{bftable}).
\begin{table}[h]
\centering
\caption{\small The scale for interpreting Bayes factors by Harold \citet{jeffreys61}.} 
\label{bftable}
\begin{tabular}{l c r}
\hline
\multicolumn{1}{l}{\textbf{$BF(M_0, M_1)$}} & \multicolumn{1}{r}{ } &\multicolumn{1}{r}{{Strength of evidence}} \\ 
\hline
$<1:1$ & \hspace{8mm} & Negative (supports $M_1$)\\
$1:1$ to $3:1$ & & Barely worth mentioning\\
$3:1$ to $10:1$ & & Substantial\\
$10:1$ to $30:1$ & & Strong\\
$30:1$ to $100:1$ & & Very strong\\
$>100:1$ & & Decisive\\
\hline
\multicolumn{3}{l}{{\scriptsize{For a detailed description of Bayes factors see \citet{kass95}}}} 
\end{tabular}
\end{table}


Unfortunately, direct calculation of the posterior odds to prior odds ratio is unfeasible for most phylogenetic models. However, we can further define the posterior odds ratio as:
\begin{align*}
\frac{\mathbb{P}(M_0 \mid \mathbf X)}{\mathbb{P}(M_1 \mid \mathbf X)} = \frac{\mathbb{P}(M_0)}{\mathbb{P}(M_1)} \frac{\mathbb{P}(\mathbf X \mid M_0)}{\mathbb{P}(\mathbf X \mid M_1)},
\end{align*}
where $\mathbb{P}(\mathbf X \mid M_i)$ is the \textit{marginal likelihood} of the data marginalized over all parameters for $M_i$; it is also referred to as the \textit{model evidence} or \textit{integrated likelihood}.
More explicitly, the marginal likelihood is the probability of the set of observed data ($\mathbf X$) under a given model ($M_i$), while averaging over all possible values of the parameters of the model ($\theta_i$) with respect to the prior density on $\theta_i$
\begin{align}\label{margeLike}
\mathbb{P}(\mathbf X \mid M_i) = \int \mathbb{P}(\mathbf X \mid \theta_i) \mathbb{P}(\theta_i)dt.
\end{align}
If you refer back to equation \ref{BFeq1}, you can see that, with very little algebra, the ratio of marginal likelihoods is equal to the Bayes factor:
\begin{align}\label{bfFormula}
BF(M_0,M_1) = \frac{\mathbb{P}(\mathbf X \mid M_0)}{\mathbb{P}(\mathbf X \mid M_1)} = \frac{\mathbb{P}(M_0 \mid \mathbf X, \theta_0)}{\mathbb{P}(M_1 \mid \mathbf X, \theta_1)} \div \frac{\mathbb{P}(M_0)}{\mathbb{P}(M_1)}. 
\end{align}
Therefore, we can perform a Bayes factor comparison of two models by calculating the marginal likelihood for each one. % Simple as pie, right?
Alas, exact solutions for calculating marginal likelihoods are not known for phylogenetic models (see equation \ref{margeLike}), thus we must resort to numerical integration methods to estimate or approximate these values. 
In this exercise, we will estimate the marginal likelihood for each partition scheme
using both the stepping-stone \citep{xie11} and path sampling estimators \citep{gelman1998,lartillot06, friel2008}. 

\bigskip
\subsection{Getting Started}

%This is an in-depth tutorial on model selection and data partitioning in RevBayes and this document contains a lot of information. We will \textbf{walk through} the unpartitioned model section (Section \ref{secUnif}) and the remaining sections will be available for you to try on your own. 

This tutorial assumes that you have already downloaded, compiled, and installed RevBayes. 
We also recommend that---if you are working on a Unix machine---you put the {\tt rb} binary in your path.

For the exercises outlined in this tutorial, we will use RevBayes interactively by typing commands in the command-line console.
The format of this exercise uses \colorbox{shadecolor}{\tt lavender blush shaded boxes} to delineate important steps. 
The various RevBayes commands and syntax are specified using \cl{typewriter text}. And the specific commands that you should type (or copy/paste) into RevBayes are indicated by shaded box and prompt. For example, after opening the RevBayes program, you can load your data file:

{\tt \begin{snugshade*}
\begin{lstlisting}
RevBayes > data_atpB <- readCharacterData("data/conifer_atpB.nex")[1]
\end{lstlisting}
\end{snugshade*}}

For this command, type in the command and its options:\\ \cl{data\_atpB <- readCharacterData("data/conifer\_atpB.nex")[1]}. \textbf{DO NOT} type in ``\cl{RevBayes >}'', the prompt is simply included to replicate what you see on your screen. 

Multi-line entries, particularly loops, will often be displayed in boxes without the \cl{RevBayes >} prompt so that they can be copied and pasted wholly. 
{\tt \begin{snugshade*}
\begin{lstlisting}
for( i in 1:12 ){
  x[i] ~ dnExponential(1.0)
}
\end{lstlisting}
\end{snugshade*}}


%Screen output resulting from the execution of a function or command will be in \colorbox{shadecolor}{\textcolor{outputcol}{\cl{forest green typewriter text in a shaded box}}}: 
%
%{\tt\begin{snugshade*}
%{\textcolor{outputcol}{ 
%\begin{lstlisting}
%Successfully read one character matrix from file `data/conifer\_dna.nex'
%\end{lstlisting}}}
%\end{snugshade*}}


This tutorial also includes hyperlinks: bibliographic citations are {\textcolor{citescol}{burnt orange}} and link to the full citation in the references, external URLs are {\textcolor{urlscol}{cerulean}}, and internal references to figures and equations are {\textcolor{linkscol}{purple}}.

The various exercises in this tutorial take you through the steps required to perform phylogenetic analyses of the example datasets. 
In addition, we have provided the output files for every exercise so you can verify your results. (Note that since the MCMC runs you perform will start from different random seeds, the output files resulting from your analyses \textit{will not} be identical to the ones we provide you.)

\exs{Download data and output files from: \href{https://molevol.mbl.edu/index.php/RevBayes}{https://molevol.mbl.edu/index.php/RevBayes}
%\href{http://treethinkers.org/phylogenetic-inference-using-mrbayes-v3-2/}{\small link}
}


\exs{Also note that ``pre-cooked'' output files are provided in the download. Throughout this tutorial, you can use those files to summarize output if you do not have time to run the full analyses yourself.}


%\exs{Open the Rev file, \taha{Rev file name}, in a text editor.  This file contains 
%all of the commands required to perform the necessary analyses to explore various
%partition schemes (unpartitioned, partitioned by gene region, and partitioned by
%gene region$+$codon position).  The details of each command are described in adjacent comments, after a \cl{\#}; {\it e.g.}, \cl{\# this is a Rev comment}. }

%Typically, we would perform these analyses by simply sourcing these files in RevBayes. For the purposes of this exercise, however, we will walk through the different steps interactively in the command line.

\bigskip
\subsection{Launch RevBayes}

Execute the RevBayes binary. If this program is in your path, then you can simply type in your Unix terminal:

\exs{\cl{\$ rb}}



\bigskip
\subsection{An Unpartitioned Analysis}\label{secUnif} 

The first model is a uniform model over all sequences.
Hence, the model is equivalent to the exercise on substitution models.
To specify the model please consult the previous exercise.


\bigskip
{\large \textcolor{mycol}{\textsc{Estimating the Marginal Likelihood}}}

Typically, model comparison is performed prior to running the full MCMC analysis under a model. 
If you calculated the Bayes factors to determine the relative support for the uniform model and found that there was strong evidence supporting this model over others (hint: this is not true if you proceed with this tutorial), then it would be worth your time to proceed with the MCMC steps outlined above. 
The following steps will describe using stepping-stone and path sampling methods on a set of power posteriors to estimate marginal likelihoods under the uniform model. 

With a fully specified model, we can set up the \cl{powerPosterior()} analysis to create a file of `powers' and likelihoods from which we can estimate the marginal likelihood using stepping-stone or path sampling. 
This method computes a vector of powers from a beta distribution, then executes an MCMC run for each power step while raising the likelihood to that power. In this implementation, the vector of powers starts with 1, sampling the likelihood close to the posterior and incrementally sampling closer and closer to the prior as the power decreases. 



%\textbf{\textit{Clear Workspace and Load the Data and Model}}

Just to be safe, it is better to clear the workspace and re-load the data and model:
{\tt \begin{snugshade*}
\begin{lstlisting}
RevBayes > clear()
RevBayes > source("RevBayes_scripts/uniform_partition_model.Rev")
\end{lstlisting}
\end{snugshade*}}


First, we create the variable containing the power posterior. 
This requires us to provide a model and vector of moves, as well as an output file name. 
The \cl{cats} argument sets the number of power steps.
{\tt \begin{snugshade*}
\begin{lstlisting}
RevBayes > pow_p <- powerPosterior(mymodel, moves, "pow_p_uniform.out", cats=50) 
\end{lstlisting}
\end{snugshade*}}

We can start the power posterior by first burning in the chain and and discarding the first 10000 states.  
{\tt \begin{snugshade*}
\begin{lstlisting}
RevBayes > pow_p.burnin(generations=10000,tuningInterval=1000)
\end{lstlisting}
\end{snugshade*}}

Now execute the run with the \cl{.run()} function:
{\tt \begin{snugshade*}
\begin{lstlisting}
RevBayes > pow_p.run(generations=1000)  
\end{lstlisting}
\end{snugshade*}}

Once the power posteriors have been saved to file, create a stepping stone sampler. This function can read any file of power posteriors and compute the marginal likelihood using stepping-stone sampling. 
{\tt \small \begin{snugshade*}
\begin{lstlisting}
RevBayes > ss <- steppingStoneSampler(file="pow_p_uniform.out", powerColumnName="power", likelihoodColumnName="likelihood")
\end{lstlisting}
\end{snugshade*}}

Compute the marginal likelihood under stepping-stone sampling using the member function \cl{marginal()} of the \cl{ss} variable and record the value in Table \ref{ssTable}.
{\tt \begin{snugshade*}
\begin{lstlisting}
RevBayes > ss.marginal() 
\end{lstlisting}
\end{snugshade*}}

Path sampling is an alternative to stepping-stone sampling and also takes the same power posteriors as input. 
{\tt \small \begin{snugshade*}
\begin{lstlisting}
RevBayes > ps <- pathSampler(file="pow_p_uniform.out", powerColumnName="power", likelihoodColumnName="likelihood")
\end{lstlisting}
\end{snugshade*}}

Compute the marginal likelihood under stepping-stone sampling using the member function \cl{marginal()} of the \cl{ps} variable and record the value in Table \ref{ssTable}.
{\tt \begin{snugshade*}
\begin{lstlisting}
RevBayes > ps.marginal() 
\end{lstlisting}
\end{snugshade*}}


%%%%%%%%

\vspace{5mm}

{\begin{framed}
\begin{center}
Stop here or continue on to evaluate partitioned models...
\end{center}
\end{framed}}



\bigskip
\subsection{Partitioning by Gene Region}\label{secByGene}

The uniform model used in the previous section assumes that all sites in the alignment evolved under the same process described by a shared tree, branch length proportions, and parameters of the GTR+$\Gamma$ substitution model.
However, our alignment contains two distinct gene regions---atpB and rbcL---so we may wish to explore the possibility that the substitution process differs between these two gene regions.
This requires that we first specify the data partitions corresponding to these two genes, then define an independent substitution model for each data partition. 

\textbf{\textit{Clear Workspace and Reload Data}}

{\tt \begin{snugshade*}
\begin{lstlisting}
RevBayes > clear()
\end{lstlisting}
\end{snugshade*}}

Since we wish to avoid individually specifying each parameter of the GTR+$\Gamma$ model for each of our data partitions, we can \textit{loop} over our datasets and create vectors of nodes.
To do this, we begin by creating a vector of data file names:
{\tt \begin{snugshade*}
\begin{lstlisting}
RevBayes > filenames <- v("data/conifer_atpB.nex", "data/conifer_rbcL.nex")
\end{lstlisting}
\end{snugshade*}}

Set a variable for the number of partitions:
{\tt \begin{snugshade*}
\begin{lstlisting}
RevBayes > n_parts <- filenames.size()
\end{lstlisting}
\end{snugshade*}}

And create a vector of data matrices called \cl{data}:
{\tt \begin{snugshade*}
\begin{lstlisting}
for (i in 1:n_parts){
   data[i] <- readCharacterData(filenames[i])[1]
}
\end{lstlisting}
\end{snugshade*}}

Next, we can initialize some important variables. This does require, however, that both of our alignments have the same number of species and matching tip names.
{\tt \begin{snugshade*}
\begin{lstlisting}
RevBayes > n_species <- data[1].ntaxa()
RevBayes > names <- data[1].names()
RevBayes > n_branches <- 2 * n_species - 3
\end{lstlisting}
\end{snugshade*}}


\textbf{\textit{Specify the Parameters by Looping Over Partitions}}

We can avoid creating unique names for every node in our model if we use a \cl{for} loop to iterate over our partitions. Thus, we will only have to type in our entire GTR+$\Gamma$ model parameters once. 
This will produce a vector for each of the unlinked parameters --- e.g., there will be a vector of \cl{shape} nodes where the stochastic node for the first partition (atpB) will be \cl{shape[1]} and the stochastic node for the second partition (rbcL) will be called \cl{shape[2]}.
{\tt \small \begin{snugshade*}
\begin{lstlisting}
mi <- 0 # an iterator for the move vector
for (i in 1:n_parts){
  ## index i=1 : atpB gene ##
  ## index i=2 : rbcL gene ##
  
  # Exchangeability rates #
  er_prior[i] <- v(1,1,1,1,1,1)
  er[i] ~ dnDirichlet(er_prior[i])
  moves[mi++] <- mvSimplexElementScale(er[i], alpha=10, tune=true, weight=3) 

  # Stationary base frequencies #
  sf_prior[i] <- v(1,1,1,1)
  sf[i] ~ dnDirichlet(sf_prior[i])
  moves[mi++] <- mvSimplexElementScale(sf[i], alpha=10, tune=true, weight=2) 

  # Instantaneous rate matrix (deterministic) #
  Q[i] := gtr(er[i],sf[i]) 

  # Gamma-dist site rates #
  shape_prior[i] <- 0.05 
  shape[i] ~ dnExponential( shape_prior[i] )
  gamma_rates[i] := discretizeGamma( shape[i], shape[i], 4 )
  moves[mi++] <- mvScale(shape[i], lambda=0.8, tune=true, weight=3.0)
}
\end{lstlisting}
\end{snugshade*}}


\textbf{\textit{Uniform Topology and Branch Lengths}}

Our two genes evolve under different GTR rate matrices with different mean-one gamma distributions on the site rates.
However, we do assume that they share a single topology and set of branch lengths.
{\tt \begin{snugshade*}
\begin{lstlisting}
# Unrooted tree topology distribution #
topology ~ dnUniformTopology(n_species, names)

# Tree topology moves #
moves[mi++] <- mvNNI(topology, weight=10.0)
moves[mi++] <- mvSPR(topology, weight=5.0)

#### Specify a prior and moves on the branch lengths #### 
# Create a vector of branch-length variables using a for loop #
for (i in 1:n_branches) {
  br_lens[i] ~ dnExponential(10.0)
  moves[mi++] <- mvScale(br_lens[i],lambda=1,tune=true,weight=1) 
}

# A deterministic node for the tree length #
tree_length := sum(br_lens)

# Build the tree by combining the topology with br_lens #
phylogeny := treeAssembly(topology, br_lens)
\end{lstlisting}
\end{snugshade*}}

\textbf{\textit{Putting it All Together}}

Since we have a rate matrix and a site-rate model for each partition, we must create a phylogenetic CTMC for each gene. 
Additionally, we must fix the values of these nodes by attaching their respective data matrices.
These two nodes are linked by the \cl{phylogeny} node and their log-likelihoods are added to get the likelihood of the whole DAG.
{\tt \begin{snugshade*}
\begin{lstlisting}
for (i in 1:n_parts){
  phyloSeq[i] ~ dnPhyloCTMC(tree=phylogeny, Q=Q[i], siteRates=gamma_rates[i], nSites=data[i].nchar(1), type="DNA")
  phyloSeq[i].clamp(data[i])
}
\end{lstlisting}
\end{snugshade*}}


And we can pass in a single, shared node to wrap up our model DAG:
{\tt \begin{snugshade*}
\begin{lstlisting}
RevBayes > mymodel <- model(topology)
\end{lstlisting}
\end{snugshade*}}

\subsubsection*{Estimating the Marginal Likelihood}

Now run the power posterior analysis on the two-gene model.
{\tt \begin{snugshade*}
\begin{lstlisting}
RevBayes > pow_p <- powerPosterior(mymodel, moves, file="pow_p_twogene.out", cats=50) 
RevBayes > pow_p.burnin(generations=1000,tuningInterval=100)
RevBayes > pow_p.run(generations=1000)  
\end{lstlisting}
\end{snugshade*}}

Calculate the marginal likelihoods under stepping-stone sampling:
{\tt \begin{snugshade*}
\begin{lstlisting}
RevBayes > ss <- steppingStoneSampler(file="pow_p_twogene.out", powerColumnName="power", likelihoodColumnName="likelihood")
RevBayes > ss.marginal() 
\end{lstlisting}
\end{snugshade*}}

And under path sampling:
{\tt \begin{snugshade*}
\begin{lstlisting}
RevBayes > ps <- pathSampler(file="pow_p_twogene.out", powerColumnName="power", likelihoodColumnName="likelihood")
RevBayes > ps.marginal() 
\end{lstlisting}
\end{snugshade*}}

Record the marginal likelihoods in Table \ref{ssTable}.



\bigskip
\subsection{Partitioning by Codon Position and by Gene}\label{secExtremeP}

Because of the genetic code, we often find that different positions within a codon (first, second, and third) evolve at different rates.
Thus, using our knowledge of biological data, we can devise a third approach that further partitions our alignment. 
For this exercise, we will partition sites within the rbcL gene by codon position.

\textbf{\textit{Clear Workspace and Reload Data}}

{\tt \begin{snugshade*}
\begin{lstlisting}
RevBayes > clear()
RevBayes > data[1] <- readCharacterData("data/conifer_atpB.nex")[1]
RevBayes > data_rbcL <- readCharacterData("data/conifer_rbcL.nex")[1]
\end{lstlisting}
\end{snugshade*}}

\textbf{\textit{Specify Data Matrices for Each Codon Position}}

We must now add our codon-partitions to the \cl{data} vector that already contains the matrix for atpB in the first index. 
Thus, the second index will be the rbcL codon position 1. 
We can create this by calling the helper function \cl{setCodonPartition()}, which is a member function of the data matrix. 
We are assuming that the gene is \textit{in frame}, meaning the first column in your alignment is a first codon position. 
The \cl{setCodonPartition()} function takes a single argument, the position of the alignment you wish to extract. 
It then returns every third column, starting at the index provided as an argument.

Before we can use the use the \cl{setCodonPartition()} function, we must first populate the position in the \cl{data} matrix with some sequences. 
Then we call the member function of \cl{data[2]} to exclude all but the 1$^{st}$ positions.
{\tt \begin{snugshade*}
\begin{lstlisting}
RevBayes > data[2] <- data_rbcL
RevBayes > data[2].setCodonPartition(1)
\end{lstlisting}
\end{snugshade*}}

Assign the 2$^{nd}$ codon positions to \cl{data[3]}:
{\tt \begin{snugshade*}
\begin{lstlisting}
RevBayes > data[3] <- data_rbcL
RevBayes > data[3].setCodonPartition(2)
\end{lstlisting}
\end{snugshade*}}

Assign the 3$^{rd}$ codon positions to \cl{data[4]}:
{\tt \begin{snugshade*}
\begin{lstlisting}
RevBayes > data[4] <- data_rbcL
RevBayes > data[4].setCodonPartition(3)
\end{lstlisting}
\end{snugshade*}}



Now we can query the vector of data matrices to get the size, which is 4:
{\tt \begin{snugshade*}
\begin{lstlisting}
RevBayes > n_parts <- data.size()
\end{lstlisting}
\end{snugshade*}}

And set the special variables from the data:
{\tt \begin{snugshade*}
\begin{lstlisting}
RevBayes > n_species <- data[1].ntaxa()
RevBayes > names <- data[1].names()
RevBayes > n_branches <- 2 * n_species - 3
\end{lstlisting}
\end{snugshade*}}




\textbf{\textit{Specify the Parameters by Looping Over Partitions}}

Setting up the GTR+$\Gamma$ model is just like in the two-gene analysis, except this time \cl{n\_parts} is equal to 4, so now our vectors of stochastic nodes should all contain nodes for each of the partitions.
{\tt \small \begin{snugshade*}
\begin{lstlisting}
mi <- 0 # an iterator for the move vector
for (i in 1:n_parts){
  ## index i=1 : atpB gene
  ## index i=2 : rbcL gene position 1
  ## index i=3 : rbcL gene position 2
  ## index i=4 : rbcL gene position 3

  # Exchangeability rates #
  er_prior[i] <- v(1,1,1,1,1,1)
  er[i] ~ dnDirichlet(er_prior[i])
  moves[mi++] <- mvSimplexElementScale(er[i], alpha=10, tune=true, weight=3) 

  # Stationary base frequencies #
  sf_prior[i] <- v(1,1,1,1)
  sf[i] ~ dnDirichlet(sf_prior[i])
  moves[mi++] <- mvSimplexElementScale(sf[i], alpha=10, tune=true, weight=2) 

  # Instantaneous rate matrix (deterministic) #
  Q[i] := gtr(er[i],sf[i]) 

  # Gamma-dist site rates #
  shape_prior[i] <- 0.05 
  shape[i] ~ dnExponential( shape_prior[i] )
  gamma_rates[i] := discretizeGamma( shape[i], shape[i], 4 )
  moves[mi++] <- mvScale(shape[i], lambda=0.8, tune=true, weight=3.0)
}
\end{lstlisting}
\end{snugshade*}}


\textbf{\textit{Uniform Topology and Branch Lengths}}

We are still assuming that the genes share a single topology and branch lengths.
{\tt \begin{snugshade*}
\begin{lstlisting}
# Unrooted tree topology distribution #
topology ~ dnUniformTopology(n_species, names)

# Tree topology moves #
moves[mi++] <- mvNNI(topology, weight=10.0)
moves[mi++] <- mvSPR(topology, weight=5.0)

#### Specify a prior and moves on the branch lengths #### 
# Create a vector of branch-length variables using a for loop #
for (i in 1:n_branches) {
  br_lens[i] ~ dnExponential(10.0)
  moves[mi++] <- mvScale(br_lens[i],lambda=1,tune=true,weight=1) 
}

# A deterministic node for the tree length #
tree_length := sum(br_lens)

# Build the tree by combining the topology with br_lens #
phylogeny := treeAssembly(topology, br_lens)
\end{lstlisting}
\end{snugshade*}}

\textbf{\textit{Putting it All Together}}

We must specify a phylogenetic CTMC node for each of our partition models.
{\tt \begin{snugshade*}
\begin{lstlisting}
for (i in 1:n_parts){
  phyloSeq[i] ~ dnPhyloCTMC(tree=phylogeny, Q=Q[i], siteRates=gamma_rates[i], nSites=data[i].nchar(1), type="DNA")
  phyloSeq[i].clamp(data[i])
}
\end{lstlisting}
\end{snugshade*}}


And then wrap up the DAG using the \cl{model()} function:
{\tt \begin{snugshade*}
\begin{lstlisting}
RevBayes > mymodel <- model(topology)
\end{lstlisting}
\end{snugshade*}}


\subsubsection*{Estimating the Marginal Likelihood}

Sample likelihoods from the set of power posteriors:
{\tt \begin{snugshade*}
\begin{lstlisting}
RevBayes > pow_p <- powerPosterior(mymodel, moves, file="pow_posterior_genecodon.out", cats=50) 
RevBayes > pow_p.burnin(generations=1000,tuningInterval=100)
RevBayes > pow_p.run(generations=1000)  
\end{lstlisting}
\end{snugshade*}}

Compute the stepping-stone estimate of the marginal likelihood:
{\tt \begin{snugshade*}
\begin{lstlisting}
RevBayes > ss <- steppingStoneSampler(file="pow_posterior_genecodon.out", powerColumnName="power", likelihoodColumnName="likelihood")
RevBayes > ss.marginal() 
\end{lstlisting}
\end{snugshade*}}

Compute the path-sampling estimate of the marginal likelihood:
{\tt \begin{snugshade*}
\begin{lstlisting}
RevBayes > ps <- pathSampler(file="pow_posterior_genecodon.out", powerColumnName="power", likelihoodColumnName="likelihood")
RevBayes > ps.marginal() 
\end{lstlisting}
\end{snugshade*}}

Now record the marginal likelihoods in Table \ref{ssTable}.



\bigskip
\subsection{Compute Bayes Factors and Select Model}


Now that we have estimates of the marginal likelihood under each of our different models, we can evaluate their relative plausibility using Bayes factors.
Use Table \ref{ssTable} to summarize the marginal log-likelihoods estimated using the stepping-stone and path-sampling methods.
\begin{Form}
\begin{table}[h]
\centering
\caption{\small Estimated marginal likelihoods for different partition configurations$^*$.}
\begin{tabular}{l c c c c}
\hline
\multicolumn{1}{l}{\textbf{ }} &\multicolumn{1}{r}{\textbf{ }} & \multicolumn{3}{c}{\textbf{Marginal lnL estimates}} \\ 
\cline{3-5}
\multicolumn{1}{l}{\textbf{Partition}} & \multicolumn{1}{r}{\hspace{3mm}} & \multicolumn{1}{c}{\textit{Stepping-stone}} & \multicolumn{1}{r}{\hspace{3mm}} & \multicolumn{1}{c}{\textit{Path sampling}} \\ 
\hline
\ref{secUnif} uniform ($M_1$) & \hspace{15mm} & \TextField[name=m1,backgroundcolor={.85 .85 .85},color={1 0 0},height=4ex]{}  & \hspace{15mm} & \TextField[name=ml2,backgroundcolor={.85 .85 .85},color={0 0 1},height=4ex]{} \\
\hline
\ref{secByGene} moderate ($M_2$) & \hspace{3mm} &\TextField[name=ml3,backgroundcolor={.85 .85 .85},color={1 0 0},height=4ex]{}   & \hspace{3mm} & \TextField[name=ml4,backgroundcolor={.85 .85 .85},color={0 0 1},height=4ex]{} \\
\hline
\ref{secExtremeP} extreme ($M_3$) & \hspace{3mm} & \TextField[name=ml5,backgroundcolor={.85 .85 .85},color={1 0 0},height=4ex]{} & \hspace{3mm} & \TextField[name=ml6,backgroundcolor={.85 .85 .85},color={0 0 1},height=4ex]{} \\
\hline
{\footnotesize{$^*$you can edit this table}}\\
\end{tabular}
\label{ssTable}
\end{table}
\end{Form}

Phylogenetics software programs log-transform the likelihood to avoid \href{http://en.wikipedia.org/wiki/Arithmetic_underflow}{underflow}, because multiplying likelihoods results in numbers that are too small to be held in computer memory.
Thus, we must use a different form of equation \ref{bfFormula} to calculate the ln-Bayes factor (we will denote this value $\mathcal{K}$):
\begin{align}\label{LNbfFormula}
\mathcal{K}=\ln[BF(M_0,M_1)] = \ln[\mathbb{P}(\mathbf X \mid M_0)]-\ln[\mathbb{P}(\mathbf X \mid M_1)],
\end{align}
where $\ln[\mathbb{P}(\mathbf X \mid M_0)]$ is the \textit{marginal lnL} estimate for model $M_0$. 
The value resulting from equation \ref{LNbfFormula} can be converted to a raw Bayes factor by simply taking the exponent of $\cal{K}$
\begin{align}\label{LNbfFormula2}
BF(M_0,M_1) = e^{\cal{K}}.
\end{align}
Alternatively, you can interpret the strength of evidence in favor of $M_0$ using the $\cal{K}$ and skip equation \ref{LNbfFormula2}. 
In this case, we evaluate the $\cal{K}$ in favor of model $M_0$ against model $M_1$ so that:
\begin{center}
\begin{tabular}{l}
if $\mathcal{K} > 1$, then model $M_0$ wins\\
if $\mathcal{K} < -1$, then model $M_1$ wins.
\end{tabular}
\end{center}
Thus, values of $\mathcal{K}$ around 0 indicate ambiguous support. 

%Alternatively, you can refer to a different scale for interpreting the $\ln[BF(M_0,M_1)]$ described by \citet{kass95}.
%\begin{table}[h]
%\centering
%\caption{\small The scale for interpreting $log_e$ Bayes factors by \citet{kass95}.} 
%\begin{tabular}{l c r}
%\hline
%\multicolumn{1}{l}{\textbf{$2\ln[BF(M_0, M_1)]$}} & \multicolumn{1}{r}{ } &\multicolumn{1}{r}{{Strength of evidence \textit{against} $M_0$}} \\ 
%\hline
%0 to 2 & \hspace{8mm} & Not worth more than a bare mention\\
%2 to 6 & \hspace{8mm} & Positive\\
%6 to 10 & \hspace{8mm} & Strong\\
%$>10$ & \hspace{8mm} & Very strong\\
%\hline
%\end{tabular}
%\label{bftable2}
%\end{table}



Using the values you entered in Table \ref{ssTable} and equation \ref{LNbfFormula},  calculate the ln-Bayes factors (using $\mathcal{K}$) for the different model comparisons. 
Enter your answers in Table \ref{bfTable} using the stepping-stone and the path-sampling estimates of the marginal log likelihoods. 

\begin{Form}
\begin{table}[h!]
\centering
\caption{\small Bayes factor calculation$^*$.}
\begin{tabular}{l c c c c}
\hline
\multicolumn{1}{l}{\textbf{ }} &\multicolumn{1}{r}{\textbf{ }} & \multicolumn{3}{c}{\textbf{ln-Bayes Factor} ($\mathcal{K}$)} \\ 
\cline{3-5}
\multicolumn{1}{l}{\textbf{Model comparison}} & \multicolumn{1}{r}{\hspace{3mm}} & \multicolumn{1}{c}{\textit{Stepping-stone}} & \multicolumn{1}{r}{\hspace{3mm}} & \multicolumn{1}{c}{\textit{Path sampling}} \\ 
\hline
$M_1,M_2$ & \hspace{15mm} & \TextField[name=ml7,backgroundcolor={.85 .85 .85},color={1 0 0},height=4ex]{}  & \hspace{15mm} & \TextField[name=ml8,backgroundcolor={.85 .85 .85},color={0 0 1},height=4ex]{} \\
\hline
$M_2,M_3$ & \hspace{3mm} & \TextField[name=ml9,backgroundcolor={.85 .85 .85},color={1 0 0},height=4ex]{} & \hspace{3mm} & \TextField[name=ml10,backgroundcolor={.85 .85 .85},color={0 0 1},height=4ex]{} \\
\hline
$M_1,M_3$ & \hspace{3mm} &  \TextField[name=ml11,backgroundcolor={.85 .85 .85},color={1 0 0},height=4ex]{} & \hspace{3mm} & \TextField[name=ml12,backgroundcolor={.85 .85 .85},color={0 0 1},height=4ex]{} \\
\hline
Supported model? & \hspace{3mm} &  \TextField[name=ml13,backgroundcolor={1 .85 .85},color={1 0 0},height=4ex]{} & \hspace{3mm} & \TextField[name=ml14,backgroundcolor={.85 .85 1},color={0 0 1},height=4ex]{} \\
\hline
{\footnotesize{$^*$you can edit this table}}\\
\end{tabular}
\label{bfTable}
\end{table}
\end{Form}

Once you complete Table \ref{bfTable}, you will notice that the Bayes factor comparison indicates strong evidence in support of the highly partitioned model using both the stepping-stone and path sampling estimates of the marginal likelihoods. 
However, this does not mean that model $M_3$ is the \textit{true} partition model.
We only considered three out of the many, many possible partitions for 2,659 sites (the number of possible partitions can be viewed if you compute the \href{http://www.wolframalpha.com/input/?i=bell+number(2659)}{$2659^{th}$ Bell number}). 
Given the strength of support for the highly partitioned model, it is possible that further partitioning is warranted for these data. 
In particular, partitioning the dataset by codon position for both atpB \textit{and} rbcL is an important next step for this exercise (consider taking some time on your own to test this model). 

Because of the computational costs of computing marginal likelihoods and the vast number of possible partitioning strategies, it is not feasible to evaluate all of them. 
New methods based on nonparametric Bayesian models have recently been applied to address this problem \citep{lartillot04,huelsenbeck07b,wuCH2013}. 
These approaches use an infinite mixture model \citep[the Dirichlet process;][]{ferguson73,antoniak74} that places non-zero probability on \textit{all} of the countably-infinite possible partitions for a set of sequences. 
Bayesian phylogenetic inference under these models is implemented in the program \href{http://megasun.bch.umontreal.ca/People/lartillot/www/index.htm}{PhyloBayes} \citep{Lartillot2009} and the \href{https://code.google.com/p/subst-bma/}{subst-bma} plug-in for \href{http://beast2.cs.auckland.ac.nz/index.php/Main_Page}{BEAST2} \citep{wuCH2013}.


% $$B_{1,2} = e^{ln(B_{12})}$$
% $$ln(MlnL_1) - ln(MlnL_2)$$

%When you complete Table \ref{bfTable} and evaluate the Bayes factors using \citet{jeffreys61} scale in Table \ref{bftable}, you will notice that no single model model stands out as ``strongly supported''. 
%Every Bayes factor comparison is ``barely worth mentioning''. 

Note that Bayes factors based on comparison of HM-based marginal likelihoods 
often \textit{strongly} favor the most extremely partitioned mixed model. 
In fact, the harmonic mean estimator has been shown to provide unreliable estimates of marginal likelihoods, compared to more robust approaches \citep{lartillot06,xie11,fan11}.
Based on these studies, it is recommended that you avoid %never, ever, ever
using HM-derived marginal likelihoods for Bayes factor comparisons.
%base Bayes factors on comparisons of HM-derived marginal likelihoods. 
(The Canadian Bayesian Radford Neal says the harmonic mean is the ``\href{http://radfordneal.wordpress.com/2008/08/17/the-harmonic-mean-of-the-likelihood-worst-monte-carlo-method-ever/}{worst Monte Carlo method ever}''.)



\bigskip
\subsection{Perform MCMC Analysis Under Preferred Model}

\textbf{\textit{Clear Workspace and Load the Data and Model}}


{\tt \begin{snugshade*}
\begin{lstlisting}
RevBayes > clear()
RevBayes > source("RevBayes_scripts/<preferred>_partition_model.Rev")
\end{lstlisting}
\end{snugshade*}}

\textbf{\textit{Specify Monitors}}

{\tt \begin{snugshade*}
\begin{lstlisting}
RevBayes > monitors[1] <- mnModel(filename="conifer_prefmodel_mcmc.log", printgen=100)
\end{lstlisting}
\end{snugshade*}}


{\tt \begin{snugshade*}
\begin{lstlisting}
RevBayes > monitors[2] <- mnFile(filename="conifer_prefmodel_mcmc.trees",printgen=100, phylogeny)
\end{lstlisting}
\end{snugshade*}}


{\tt \begin{snugshade*}
\begin{lstlisting}
RevBayes > monitors[3] <- mnScreen(printgen=10, separator = " | ", tree_length)
\end{lstlisting}
\end{snugshade*}}

\textbf{\textit{Initialize and Run MCMC}}

{\tt \begin{snugshade*}
\begin{lstlisting}
RevBayes > mymcmc <- mcmc(mymodel, monitors, moves)
\end{lstlisting}
\end{snugshade*}}


{\tt \begin{snugshade*}
\begin{lstlisting}
RevBayes > mymcmc.burnin(generations=10000,tuningInterval=1000)
\end{lstlisting}
\end{snugshade*}}


{\tt \begin{snugshade*}
\begin{lstlisting}
RevBayes > mymcmc.run(generations=30000)
\end{lstlisting}
\end{snugshade*}}


\bigskip
\subsection{Summarize and Analyze MCMC Output}

{\tt \begin{snugshade*}
\begin{lstlisting}
RevBayes > treetrace <- readTreeTrace("conifer_prefmodel_mcmc.trees")
RevBayes > treetrace.summarize()
\end{lstlisting}
\end{snugshade*}}


{\tt \begin{snugshade*}
\begin{lstlisting}
RevBayes > mapTree(treetrace,"conifer_prefmodel_MAP.tre")
\end{lstlisting}
\end{snugshade*}}




The trees in these files are also annotated with various branch- or node-specific parameters or statistics in an extended Newick format called NHX. 
We can use FigTree to visualize these summary trees. 
\exs{Open the summary tree in FigTree: \cl{conifer\_prefmodel\_MAP.tre}.}
\exs{Use the tools on the side panel to display the posterior probabilities as node labels. }

%\begin{figure}[h]
%\centering
%\fbox{\includegraphics[width=4in,angle=0]{figures/tree_uniform_analysis.eps}}
%\caption{\small The summary tree from the uniform analysis, with posterior probabilities labeled at nodes. \taha{need to make a new figure}}
%\label{uniformFigTree}
%\end{figure}

\bigskip
\subsection*{Batch Mode}

If you wish to run this exercise in batch mode, the files are provided for you. 

You can carry out these batch commands by providing the file name when you execute the \cl{rb} binary in your unix terminal (this will overwrite all of your existing run files).
\exs{\cl{\$ rb full\_analysis.Rev}}

\bigskip
\section*{Useful Links}

\begin{itemize}
\item RevBayes: \href{https://github.com/revbayes/code}{https://github.com/revbayes/code} \\ \vspace{-7mm}
\item MrBayes: \href{http://mrbayes.sourceforge.net/}{http://mrbayes.sourceforge.net} \\ \vspace{-7mm}
\item PhyloBayes: \href{http://megasun.bch.umontreal.ca/People/lartillot/www/index.htm}{http://www.phylobayes.org} \\ \vspace{-7mm}
\item Bali-Phy: \href{http://www.bali-phy.org/}{http://www.bali-phy.org} \\ \vspace{-7mm}
%\item BEAGLE: \href{http://code.google.com/p/beagle-lib/}{http://code.google.com/p/beagle-lib/} \\ \vspace{-7mm}
\item Tree Thinkers: \href{http://treethinkers.org/}{http://treethinkers.org} \\ \vspace{-7mm}
\end{itemize}

Questions about this tutorial can be directed to: \\\vspace{-10mm}
\begin{itemize}
\item Tracy Heath (email: \href{mailto:tracyh@berkeley.edu}{tracyh@berkeley.edu}) \\\vspace{-8mm}
\item Michael Landis (email: \href{mailto:mlandis@berkeley.edu}{mlandis@berkeley.edu}) \\\vspace{-8mm} 
\item Sebastian H\"{o}hna (email: \href{mailto:sebastian.hoehna@gmail.com}{sebastian.hoehna@gmail.com})
\end{itemize}

\href{http://creativecommons.org/licenses/by/4.0/}{\includegraphics[scale=0.8]{figures/ccby.eps}} This tutorial was written by \href{http://phylo.bio.ku.edu/content/tracy-heath}{Tracy Heath}, \href{http://michaeljameslandis.com/}{Michael Landis}, and Sebastian H\"{o}hna; licensed under a \href{http://creativecommons.org/licenses/by/4.0/}{Creative Commons Attribution 4.0 International License}. 
(This tutorial is based on the \href{http://treethinkers.org/tutorials/phylogenetic-inference-using-mrbayes-v3-2/}{Phylogenetic Inference using MrBayes v3.2} tutorial written by Tracy Heath, Conor Meehan, and Brian Moore and some content is reproduced here.
Mark Holder, Ben Liebeskind, Emily McTavish, and April Wright provided helpful comments.)

Version dated: \today

\nocite{sukumaran10,rannala96,robert02,rodrigue08,ronquist09,Rubinstein1981,simon01,Smith1993,Suchard2005,verdinelli95,wong08,Yang2007,Yang2005,yang97b}
%\newpage

%%%%%%%%%%%%%%%%%%%%%%%%%%%%%%%%%%%%
%  REFERENCES  
%%%%%%%%%%%%%%%%%%%%%%%%%%%%%%%%%%%%

{\renewcommand{\bibsection}{\section*{Relevant References}}


\bibliography{bib_tex/master_refs}




\end{document}
